\section{Implementation}
\label{sec:implementation}

\subsection{Chaincode and transaction flow}

The chaincode is written in Go 1.21 against \texttt{fabric-contract-api-go} and targets Fabric 2.5. Eight transactions are exposed: \texttt{AssignRole}, \texttt{RevokeRole}, \texttt{RenewRole}, \texttt{CheckAccess}, \texttt{IssueCrossZoneToken}, \texttt{RevokeCrossZoneToken}, \texttt{GetAuditLog} and \texttt{PruneNonces}. An authorization denial is a successful chaincode response carrying \texttt{DENY} plus an audit write, so it is ordered and committed like any other transaction, and only an endorsement-level failure stops before the orderer.

Three implementation decisions shape the measurements. Every state-mutating call carries a UUID nonce recorded in world state, rejected on reuse before any state is read, which defeats replay inside the timestamp window. Every access attempt, granted or denied, writes an audit entry holding the caller's certificate hash, operation, resource, zone, decision and reason, so the audit trail is a byproduct of enforcement rather than a parallel log that can drift from it. And \texttt{CheckAccess} revokes an expired assignment in the same transaction that denies the request, which makes even a read-intent check state-mutating and therefore endorsed rather than served as a query.

Peer configuration for the ARM64 gateways uses 50 transactions per block, a 500\,ms batch timeout, a 256\,MB state cache, endorsement restricted to zone-local peers, Snappy compression on LevelDB and 5-minute gRPC keepalive, targeting the write amplification that dominates latency on microSD-backed storage, within the parameter space Abang et al.\ \citep{abang2024latency} identify as governing Fabric latency.

\subsection{Sensor firmware and residue-coded transmission}
\label{sec:firmware}

The sensor certificate lives in battery-backed RTC slow memory so it survives deep sleep, guarded by a CRC and a magic sentinel checked on each wake; a mismatch forces re-enrollment rather than transmission under an unverifiable identity. Signing uses wolfSSL's Ed25519 implementation (\texttt{esp32/main/signing.c}): the firmware hashes the payload with hardware SHA-256 and signs the resulting 32-byte digest, a prehash construction that is neither standard Ed25519 (which hashes the message internally with SHA-512) nor Ed25519ph (which prehashes with SHA-512 under a distinct domain separator). Interoperability depends on reproducing this nonstandard construction; an implementation of standard Ed25519 or Ed25519ph would not verify the historical signatures. The main loop wakes every 30 minutes, verifies the cached certificate, samples the sensors, signs the payload, transmits and returns to deep sleep.

Readings are transmitted as residues rather than plaintext values: a value is reduced modulo the pairwise-coprime set $\{\CrtModuli\}$ (\texttt{esp32/main/crt\_encode.h}, verified against the constants compiled into the firmware) and the three residues go out as separate single-byte packets, any two of which reconstruct the original by the Chinese Remainder Theorem \citep{garner1959residue}. The admissible range is where this scheme goes wrong, and the deployed firmware got it wrong: recovery from an arbitrary pair is unique only below the smallest pairwise product of the moduli ($\CrtMinBound$), not below their product ($97 \times 101 \times 103 = 1{,}009{,}591$); Section~\ref{sec:crt-audit} recomputes the consequence from the released engineering-unit columns rather than from the firmware's theoretical ceiling. A specified correction uses $\{\CrtCorrectedModuli\}$, a pairwise bound of $\CrtCorrectedBound$, and quantises soil moisture from twelve to eight bits so the packed value's legal domain reaches $\CrtCorrectedLegalMax$ and stays inside the bound; this correction is arithmetically checked in Section~\ref{sec:crt-audit} but, as Section~\ref{sec:corrected-absent} states, is not present as running firmware in this repository.

The firmware transmits the raw 12-bit soil ADC count and the raw centi-degree temperature reading (\texttt{esp32/main/main.c}); it does not itself compute an engineering-unit percentage. No file in this repository -- firmware, gateway or backend -- defines the calibration constant that would convert the released \texttt{soil\_moisture} percentage column back to that raw ADC count. Section~\ref{sec:crt-audit} states this gap explicitly and shows the headline finding is robust to it.

\subsection{Gateway services and revocation propagation}

Each gateway runs Fabric 2.5 from official ARM64 binaries. A Python enrollment service bound to the local interface accepts CSRs only from sensors in its own zone and rate-limits per source address, and a policy cache with a 301-second nominal TTL absorbs repeated checks from the same sensor during a transmission burst. Revocation publication and cache invalidation is the mechanism whose behaviour we most wanted to observe over a long window, and one branch is the whole story of Section~\ref{sec:revocation-results}: for a reachable gateway the cache is invalidated by the CRL event itself, so staleness is bounded by gossip propagation and the nominal TTL never binds, while for an unreachable one the bound falls back to the scheduled refresh interval. In the intended design, gateway reachability determines whether event-driven invalidation is available or whether the scheduled refresh path is required; because the historical trace lacks complete episode identifiers and synchronised connectivity state, it cannot establish which branch was followed in any given observed case.

The gateway's signature-verification routine (\texttt{gateway/gateway.py}, function \texttt{verify\_signature}) returns \texttt{True} immediately when the incoming reading carries no signature, and separately builds the bytes it verifies from an ASCII rendering of the decoded fields (\texttt{device\_id:reading\_id:value:zone}) rather than the packed binary struct the firmware signs. Both properties are read directly from the released source, not inferred from the trace; Section~\ref{sec:sig-audit} reports the consequence.

\subsection{Formal specification}

A three-role subset of the model is specified in TLA+ and checked with TLC \citep{lamport2002specifying,newcombe2015how}, finding no violation of the role-permission invariant under weak fairness. This is a sanity check on a small model, not verification of the deployed framework: it covers three of seven roles, omits zone and temporal predicates, and nothing connects it to the Go implementation, the binding gap given that the implemented hierarchy differs from the specified one (Section~\ref{sec:audit-policy}). The specification is in the artifact (\texttt{tla/AccessControl.tla}).

\subsection{The corrected implementation is specified, not released}
\label{sec:corrected-absent}

An earlier internal draft of this analysis referred to a separately tagged \texttt{chaincode/hrbac-corrected/} package, corrected firmware and a corrected gateway verifier, together with regression tests and a byte-for-byte signature test vector. Repository inspection performed for this version (Section~\ref{sec:methodology}, and \texttt{analysis/DATA\_DISCREPANCIES.md}) found none of these in the codebase snapshot analyzed here: there is no \texttt{hrbac-corrected} directory, no corrected \texttt{esp32/main} variant, no \texttt{policy-requirements.json} oracle file, and no \texttt{v09-deployed-historical} or \texttt{v09-corrected} tag in the repository's git history, which at the time of this analysis contains a single commit. We therefore report the corrections that Sections~\ref{sec:role-hierarchy}, \ref{sec:firmware} and \ref{sec:audit} describe as \emph{specified and arithmetically checked} -- we verified independently that the proposed moduli, hierarchy edits and verification logic satisfy the stated requirement -- rather than as \emph{implemented, tested or released}. No performance, correctness or security claim in this article is made about a corrected chaincode, firmware or gateway build, because no such build exists in the inventoried repository. This is treated as a limitation (L6, Section~\ref{sec:limitations}) rather than smoothed over.
